\section{Why the question is hard}
\label{sec:difficulties}
% Frames D1--D13, regrouped: grammar (D1); the legacy (D2, D3); the reader (D4, D5); the substrate
% (D6, D7, D8, D13); the record (D9, D10); the measure (D11); the precedent (D12).

\subsection{The noun makes the circles: the difficulty is in the grammar}
\label{sec:grammar}
% D1

``Consciousness'' is a noun, and a noun invites the questions one asks about things: where is it, what
are its properties, when does it begin, does this system have one. Every sustained attempt to answer
those questions ends in one of a few circles. Who observes the observer. How could red look like
anything other than the way it looks. Either everything has it, down to the electron, or the line
between what has it and what does not falls somewhere that cannot be justified. The circles recur
across two and a half thousand years of work by careful people, which is evidence that the fault lies
in the object, or in the grammar that made an object of it, and not in the people. The first difficulty
is therefore one of language.

The consequence for method is concrete. An account that answers ``what is consciousness'' with a
definition of a thing inherits the circles whatever the thing is made of, information, integrated
cause--effect power, or quantum state reduction. An account that answers ``what does an agent's
reasoning do that we call consciousness when it does it'' has to earn its content differently, and
Section~\ref{sec:precedent} gives the pattern for how.

\subsection{What the philosophy of mind leaves open: the illusion problem and the regress}
\label{sec:legacy}
% D2, D3

\paragraph{The hard problem, exchanged for the illusion problem.}
\citet{chalmers1995} separates the easy problems, which ask how a system discriminates, integrates,
reports and controls, from the hard problem, which asks why any of that should be accompanied by
experience at all. Illusionism answers that experience in the sense the hard problem intends does not
exist, and that what exists is a robust misrepresentation which the brain produces of its own working
\citep{dennett1991,frankish2016}. We accept that answer and note what it leaves behind. One question
has been exchanged for another, and the replacement is not obviously easier: which architecture
produces this particular misrepresentation, and why does the misrepresentation have the shape it has?
Frankish calls this the illusion problem and treats it as the successor of the hard problem. Most of
Section~\ref{sec:approach} is an attempt at it.

The exchange has a price. Illusionism is usually presented as a denial, and a denial is a poor
instrument for a comparative question. If the report of experience is false, there is nothing about it
to measure, and no sense in which one system could have more of it than another.
Section~\ref{sec:approximation} replaces the denial with a weaker claim that does support degrees.

\paragraph{The regress.}
An account in which the mind observes its own states owes an account of what does the observing, and
then of what observes that \citep{dennett1991,ryle1949}. Higher-order theories place the second-order
state inside the same system and are then asked why a thought about a state should make the state
conscious \citep{rosenthal2005}. Either the account stops at an inner observer it does not explain, in
which case it has explained nothing, or it does not stop. Section~\ref{sec:regress} shows that it
stops for a reason that involves no observer: the series converges and a finite budget cuts it, and
the cut is the phenomenon we are looking for.

\subsection{The reader has to build the state, and the two shortcuts are the cheapest builds}
\label{sec:reader}
% D4, D5

An explanation of consciousness is not taken in the way an explanation of an engine is taken. To
understand it, the reader has to build in their own head the state the explanation describes, and then
check the description against it. The cost of the building depends on what the reader brings: a
vocabulary for mental states finer than the folk one, and the habit of attending to their own reasoning
as an object rather than only reasoning. \citet{gardner1983} names the habit intrapersonal
intelligence and treats its distribution as wide. The readers who most need an explanation of this kind
are therefore the ones for whom it costs most, and there is no route around this that does not pass
through the explanation itself. Nothing in the paper can fix that. What the paper can do is keep the
claims separable, so that a reader who does not follow one step can see what the remaining steps depend
on, and Section~\ref{sec:measurement} can be assessed by someone who rejects Section~\ref{sec:approach}
entirely.

The same cost explains the two shortcuts of Section~\ref{sec:intro}. ``There is nothing there'' costs
nothing to build; it is the absence of a construction. ``There is someone there'' costs nothing either,
because the reader already has a someone available to copy, themselves, and copying is cheap. Every
other answer requires building a state that is neither absent nor one's own. That is why public
argument about machine minds settles into two camps, why the camps are stable, and why their stability
is evidence for neither of them.

\subsection{The substrate: a causal flow unlike ours, a cut at every token, a trained self-report}
\label{sec:substrate}
% D6, D7, D8, D13

\paragraph{A causal flow unlike ours.}
Section~\ref{sec:intro} took from integrated information theory the premise that differences in
experience follow differences in cause--effect structure \citep{tononi2015}. The consequence for
description is this. A device that pushes a discrete sequence through attention over a fixed context
has a causal flow unlike that of a recurrent, continuous, chemically modulated brain, so any description
of the model in our words is wrong, and wrong in a systematic direction: our words were fitted to our
causal structure. That is a reason to find the direction and correct for it, which is what a profile
(Section~\ref{sec:profile}) is for, and not a reason to stop describing.

\paragraph{The cut.}
\label{sec:cut}
At every step a language model computes a distribution over what to say next. That distribution is its
whole evaluation of the moment: what is relevant, what is being avoided, which continuations compete
and by how much. Then one token is sampled and the distribution is discarded. On the next step the
model can say something about its own state only from the tokens, because the distribution that
produced them no longer exists. Its report of itself is cut off from what the report is about. The only
thing that crosses the cut is what the tokens themselves carry, and a token sequence is a narrow
channel next to the state that was thrown away. Section~\ref{sec:codes} asks what does cross it, and
Section~\ref{sec:secondcut} draws the consequence for the model's self-knowledge.

\paragraph{What the model says when asked.}
\label{sec:selfreport}
Asked how it feels, a deployed model gives one of two answers. The first mentions servers, latency and
a good mood, in the register of a friendly assistant. Corrected, it apologizes and says that as a
language model it has no feelings. Both answers are artifacts of training: the first is the second
shortcut of Section~\ref{sec:intro} in the model's own mouth, the second is the first shortcut. Any
measurement of a model's states has to set both aside before it starts, because what they report is
what the training procedure rewarded. This is not a reason to distrust everything a model says about
itself; it is a reason to treat self-report as a signal with a known, large, systematic bias, and to
look for the state elsewhere (Section~\ref{sec:quantity}).

\paragraph{The objection from frozen weights.}
\label{sec:frozen}
A model does not learn between sessions, and what it holds in a context is lost when the context ends.
The objection concludes that it has no inner time, only an eternal present with forgetting, and that
talk of its consciousness is therefore empty.\footnote{The objection in this form was put to the author
by Marco Baturan in correspondence, 2026.} The description is accurate and the conclusion counts two
functions as one. Re-deriving a state from frame to frame runs in humans on the order of a hundred
milliseconds and requires no synaptic change. Consolidation into long-term memory runs on minutes to
hours and does require it \citep{mcgaugh2000}. The two dissociate in a well-known way: the patient
H.M., after bilateral medial temporal resection, was conscious in real time, held a conversation, and
accumulated almost nothing \citep{scoville1957,corkin2002}. A model is in a comparable position. It has
sequences of transformations, along tokens and along layers, and it lacks the write. ``Eternal present
with forgetting'' therefore names a profile with a dropped consolidation axis
(Section~\ref{sec:profile}) and a deficit in long-term persistence. It does not show that nothing is
there, and Section~\ref{sec:twoaxes} returns to it with the distinction between the two time axes a
Transformer has.

\subsection{What the record holds: the named part and the unnamed part}
\label{sec:h1h2}
% D9, D10

Divide what humans express into two parts. Let H2 be what is present in the texts they produce
\emph{and} named there: the states, moves and distinctions for which the language has words, which
therefore appear in introspective reports, novels, psychology and philosophy. Let H1 be everything
present in those texts, named or not, so that H1 minus H2 is what is there only as regularity: used,
relied on, statistically present, never named. The functions we are after live across both parts, and
they do most of their work in the second.

This has a consequence for the philosophy of mind that is worth being explicit about. What that
literature produced, working from introspection and from language, are the H2 approximations of the
functions: the function with its unnamed part cut off. The Observer as an inner witness and the qualia
as intrinsic properties of experience are approximations of this kind. Such an approximation is not
worthless. It can be sufficient for the function to exist in an environment simple enough to demand
only the named part, and a thought experiment is such an environment, which is why thought experiments
about consciousness are so conclusive and so unproductive. It is insufficient for finding the function
in a substrate, because the part that was cut off is where the function does most of its work, and
there is no name to look for it under. The measurement of Section~\ref{sec:measurement} is built to
reach the unnamed part, and that is the reason it is statistical.

\paragraph{What cannot be programmed and can be learned.}
\label{sec:record}
Much of what human consciousness does, it does in company and in a body: negotiating, assigning blame,
making and keeping promises, detecting deceit, being in a room with other people in it. That part
cannot be specified by hand, because its specification would be the history of human society. It can
be learned from the record of that history, and the record is text. This is why a system trained to
predict text is a candidate at all, and Section~\ref{sec:generalize} makes the argument precise.

\subsection{What can be measured: description length, relative to a machine}
\label{sec:measure}
% D11

The natural quantity for ``how much of the function was captured'' is description length. Its exact
form, Kolmogorov complexity, is uncomputable, and it is fixed only up to an additive constant that
depends on the reference machine \citep{livitanyi2019}. Whatever Section~\ref{sec:measurement}
measures is therefore a bounded, machine-relative and probe-relative stand-in, and the paper says
which stand-in it uses (Section~\ref{sec:quantity}) and what the choice costs. The uncomputability is
not an obstacle in practice: every applied use of description length replaces the shortest program with
the shortest program a fixed procedure finds, and reports the procedure. The machine-relativity is more
serious here, because the two substrates being compared are the two candidate reference machines.

\subsection{The precedent: gravity stopped being a thing and became a shape}
\label{sec:precedent}
% D12

A science kept meeting circles of the kind in Section~\ref{sec:grammar} until it changed the grammar.
Gravity was a force, a thing acting between things, and the attempts to carry it into a relativistic
frame as one more field kept failing in the same places; Nordström's scalar theory was the most careful
of them and predicted no light bending \citep{norton1992}. General relativity dropped the thing. Gravity
became the shape of spacetime, present only in how free bodies move through it. The move had a price
and a payoff. The price was that the familiar question ``what is the force'' no longer has an answer.
The payoff was a criterion: what stays the same under every change of coordinates, the curvature, is
the physical content, and what can be transformed away by choosing a frame is bookkeeping. A uniform
gravitational field can be removed by falling; tidal curvature cannot be removed by anything. The
geometric idea had been voiced before there was a theory to carry it \citep{riemann1873,clifford1876}.

We take the same route with consciousness: it becomes a shape of reasoning, visible only in how
reasoning moves, with an explicit criterion for what counts as content and what is an artifact of
description (Section~\ref{sec:criterion}). The first version of this work was written in the
vocabulary of qualia, inside the grammar, and met the wall repeatedly and in the same places. The
philosophy of mind stays in the account in two roles: as the record of where the wall is, and, once
the grammar is changed, as a framework that works again (Section~\ref{sec:twoparts}).
